%==============================================================================
\section{Data Distribution --- Phân phối dữ liệu}
\label{sec:data-distribution}
%==============================================================================

\begin{dinhnghia}[title={Phân phối dữ liệu}]
Trong ML, phân phối dữ liệu là cách các điểm dữ liệu \textbf{trải rộng} trên dataset.
Hiểu phân phối ảnh hưởng lớn đến hiệu năng thuật toán.
\end{dinhnghia}

\begin{ynghia}[title={Đại lượng mô tả}]
Mean, median, mode, độ lệch chuẩn, phương sai mô tả xu hướng trung tâm, độ trải và hình dạng.
\end{ynghia}

\subsection{Phân phối chuẩn (Normal / Gaussian)}

\begin{dinhnghia}[title={Normal distribution}]
Phân phối xác suất liên tục dạng \textbf{chuông}, đối xứng quanh mean; hai tham số: mean $\mu$ và độ lệch chuẩn $\sigma$.
\end{dinhnghia}

\begin{phuongphap}[title={Ứng dụng trong ML}]
Mô hình hoá sai số trong hồi quy tuyến tính; cơ sở kiểm định giả thuyết và khoảng tin cậy.
\end{phuongphap}

\begin{ynghia}[title={Quy tắc 68--95--99.7}]
Khoảng $[\mu-\sigma, \mu+\sigma]$ chứa $\approx 68\%$ quan sát; $\pm 2\sigma$: $\approx 95\%$; $\pm 3\sigma$: $\approx 99{,}7\%$.
\end{ynghia}

\begin{vidu}[title={Sinh mẫu và vẽ PDF (scipy.stats)}]
\begin{lstlisting}
import numpy as np
from scipy.stats import norm
import matplotlib.pyplot as plt

mu = 0
sigma = 1
sample = np.random.normal(mu, sigma, 1000)

x = np.linspace(mu - 3*sigma, mu + 3*sigma, 100)
pdf = norm.pdf(x, mu, sigma)

plt.figure(figsize=(7.5, 3.5))
plt.hist(sample, bins=30, density=True, alpha=0.5)
plt.plot(x, pdf)
plt.show()
\end{lstlisting}

\textbf{Output}:
\tsimg{05_data_distribution/img01.jpg}
\end{vidu}

\subsection{Phân phối lệch (Skewed)}

\begin{ynghia}[title={Skewed distribution}]
Dataset không đối xứng quanh mean; phần lớn điểm dồn một phía, đuôi kéo dài phía kia.
\begin{itemize}
  \item \textbf{Lệch trái} (negative skew): đuôi dài bên trái.
  \item \textbf{Lệch phải} (positive skew): đuôi dài bên phải.
\end{itemize}
Dữ liệu lệch có thể làm méo mô hình --- cần chuẩn hoá hoặc biến đổi.
\end{ynghia}

\begin{vidu}[title={Gamma ngẫu nhiên}]
\begin{lstlisting}
import numpy as np
import matplotlib.pyplot as plt

data = np.random.gamma(2, 1, 1000)

plt.figure(figsize=(7.5, 3.5))
plt.hist(data, bins=30)
plt.xlabel('Value')
plt.ylabel('Frequency')
plt.title('Skewed Distribution')
plt.show()
\end{lstlisting}

\textbf{Output}:
\tsimg{05_data_distribution/img02.jpg}
\end{vidu}

\subsection{Phân phối đều (Uniform)}

\begin{dinhnghia}[title={Uniform distribution}]
Mọi kết quả có \textbf{xác suất bằng nhau}; không có cụm quanh một giá trị.
Hữu ích làm chuẩn so sánh hoặc sinh số ngẫu nhiên không thiên lệch.
\end{dinhnghia}

\begin{ynghia}[title={PDF liên tục}]
\[
f(x) = \begin{cases}
1 & \text{nếu } a \le x \le b \\
0 & \text{ngược lại}
\end{cases}
\]
Mean: $\frac{a+b}{2}$; variance: $\frac{(b-a)^2}{12}$.
\end{ynghia}

\begin{vidu}[title={Uniform trên $[0,1]$}]
\begin{lstlisting}
import numpy as np
import matplotlib.pyplot as plt

uniform_data = np.random.uniform(low=0, high=1, size=10000)

plt.figure(figsize=(7.5, 3.5))
plt.hist(uniform_data, bins=50, density=True)
plt.xlabel('Value')
plt.ylabel('Frequency')
plt.title('Uniform Distribution')
plt.show()
\end{lstlisting}

\textbf{Output}:
\tsimg{05_data_distribution/img03.jpg}
\end{vidu}

\subsection{Phân phối hai đỉnh (Bimodal)}

\begin{dinhnghia}[title={Bimodal distribution}]
Phân phối có \textbf{hai mode} (hai đỉnh), thường tách bởi vùng thấp (thung lũng) ở giữa --- có thể phản ánh hai tiểu quần thể.
\end{dinhnghia}

\begin{vidu}[title={Hai cụm Gaussian}]
\begin{lstlisting}
import numpy as np
import matplotlib.pyplot as plt

bimodal_data = np.concatenate((
    np.random.normal(loc=-2, scale=1, size=5000),
    np.random.normal(loc=2, scale=1, size=5000)
))

plt.figure(figsize=(7.5, 3.5))
plt.hist(bimodal_data, bins=50, density=True)
plt.xlabel('Value')
plt.ylabel('Frequency')
plt.title('Bimodal Distribution')
plt.show()
\end{lstlisting}

\textbf{Output}:
\tsimg{05_data_distribution/img04.jpg}
\end{vidu}
